\documentclass[11pt]{article}
\usepackage[margin=1in]{geometry}
\usepackage{amsmath, amssymb}
\usepackage{fontspec}
\setmonofont{Menlo}[Scale=0.85]

% \code renders inline code in the mono font, and (via url) is allowed to break
% at path separators `/` and `.`, so long file paths wrap instead of overrunning
% the right margin.
\usepackage{xurl}
\DeclareUrlCommand\code{\urlstyle{tt}}

\title{Scott Domains in Lean / Mathlib}
\author{ScottLean4 \textperiodcentered\ for Dana Scott}
\date{}

\begin{document}
\maketitle

Does Lean's Mathlib formalize Scott domains? The honest picture, from the
\code{v4.32.2} tree: the \emph{foundations} are present (and carry Dana Scott's
name in the source), but the \emph{signature constructions} of domain theory as a
model of computation are not.

\section*{What Mathlib does have (the foundations)}
\begin{itemize}
  \item \textbf{Scott continuity} --- \code{Order/ScottContinuity.lean} (with
    \code{Prod} and \code{Complete} companions): a function preserving directed
    suprema. The file's docstring reads ``Scott continuous (referring to Dana
    Scott).''
  \item \textbf{The Scott topology} --- \code{Topology/Order/ScottTopology.lean}
    (plus the \textbf{Lawson topology}, \code{LawsonTopology.lean}, and
    Scott--Hausdorff separation).
  \item \textbf{$\omega$-complete partial orders} ---
    \code{Order/OmegaCompletePartialOrder.lean} and
    \code{Topology/OmegaCompletePartialOrder.lean}: $\omega$CPOs and
    $\omega$-continuous maps, the directed-complete setting for denotational
    semantics (referenced across $\sim$13 files). Also
    \code{Order/CompletePartialOrder.lean}.
  \item \textbf{Fixed-point theory} --- \code{Order/FixedPoints.lean}
    (Knaster--Tarski least/greatest fixed points, \code{OrderHom.lfp}/\code{gfp})
    and \code{Order/BourbakiWitt.lean}. A \code{LawfulFix} operator gives the
    Kleene least fixed point $\bigsqcup_n f^n(\bot)$ of an $\omega$-continuous
    map ($\sim$22 files) --- Scott's account of recursion, used to define
    partial/recursive functions on the \code{Part} monad.
\end{itemize}

So the core apparatus --- Scott continuity, the Scott topology, cpos, and
least-fixed-point recursion --- is genuinely formalized.

\section*{What it does not have}
\begin{itemize}
  \item \textbf{No named ``Scott domain.''} There is no
    \code{DirectedCompletePartialOrder}/\code{DCPO} class (it works through
    \code{OmegaCompletePartialOrder}), and no development of bounded-complete
    algebraic dcpos or Scott information systems.
  \item \textbf{No continuous-lattice theory} under that name (the Lawson
    topology is present, but not the continuous-lattices development).
  \item \textbf{The big absence: no $D_\infty$.} The reflexive object
    $D \cong [D \to D]$ --- the model of the untyped $\lambda$-calculus --- is
    not in Mathlib, nor is the category of domains packaged as a cartesian-closed
    category for denotational semantics. The inverse-limit solution of recursive
    domain equations has not been formalized.
\end{itemize}

\section*{The upshot}
Mathlib has the \emph{ingredients} (Scott continuity, Scott topology,
$\omega$CPOs, least fixed points) but not the \emph{signature constructions} of
domain theory as a model of computation ($D_\infty$, the cartesian-closed
category of domains, continuous lattices). That gap is an opportunity:
formalizing $D_\infty$ in Lean would be a genuinely new --- and fitting ---
contribution.

\end{document}
